\chapter{Threat Modeling Blockchain Systems}

\section{Purpose}

The first five chapters gave us the primitives:

\begin{itemize}
\tightlist
\item
  State machines.
\item
  Assets and authority.
\item
  Transactions and state transitions.
\item
  Commitments and proofs.
\item
  Keys, signatures, and authorization.
\end{itemize}

This chapter turns those primitives into a working security-review
method.

\begin{quote}
A threat model is a structured explanation of what the system is, what
can go wrong, why those failures matter, and what will be done about
them.
\end{quote}

That sounds simple. In practice, most weak threat models fail in one of
three ways:

\begin{itemize}
\tightlist
\item
  They are too vague: ``attackers may steal funds.''
\item
  They are too generic: ``check access control, reentrancy, and
  arithmetic.''
\item
  They are too decorative: a diagram exists, but it does not change what
  anyone reviews, tests, monitors, or fixes.
\end{itemize}

A useful blockchain threat model must be more concrete. It should tell a
reviewer:

\begin{lstlisting}
These are the assets.
These are the state transitions that affect them.
These actors can influence those transitions.
These assumptions must hold.
These assumptions are most fragile.
These threats deserve priority.
These invariants should be written.
These tests, proofs, mitigations, and monitoring controls should exist.
\end{lstlisting}

OWASP describes threat modeling as identifying, communicating, and
understanding threats and mitigations in the context of protecting
something valuable.\footnote{OWASP, ``Threat Modeling'',
  https://owasp.org/www-community/Threat\_Modeling.} The Threat Modeling
Manifesto frames threat modeling around four practical questions: what
are we working on, what can go wrong, what will we do about it, and did
we do a good enough job.\footnote{Threat Modeling Manifesto,
  https://www.threatmodelingmanifesto.org/.}

For blockchain systems, we will adapt those questions:

\begin{lstlisting}
What state machine are we reviewing?
Which assets and authorities matter?
Which adversarial transitions can break safety, liveness, solvency, or authorization?
Which assumptions are we relying on?
Which mitigations, invariants, tests, and monitoring controls make those assumptions defensible?
\end{lstlisting}

Threat modeling is not separate from auditing. It is how you decide
where serious review begins. OWASP's smart-contract security
verification material makes this operational: security architecture
review should identify threats, assess risks, and define mitigations for
smart-contract systems.\footnote{OWASP Smart Contract Security,
  ``SCSVS-ARCH-3: Threat Modeling'',
  https://scs.owasp.org/SCSVS/controls/SCSVS-ARCH-3/.}

\section{Learning Objectives}

By the end of this chapter, you should be able to:

\begin{itemize}
\tightlist
\item
  Turn a whitepaper, docs, codebase, deployment map, or protocol
  description into a security model.
\item
  Identify actors, assets, entry points, privileged roles, trust
  boundaries, and external dependencies.
\item
  Separate users, attackers, rational participants, Byzantine actors,
  insiders, and trusted operators.
\item
  Write explicit trust assumptions instead of vague claims about
  decentralization or security.
\item
  Map attack surfaces across chain, contract, protocol, economic,
  infrastructure, and human layers.
\item
  Use threat modeling to produce review priorities, invariants, tests,
  mitigations, and monitoring requirements.
\item
  Recognize threat-modeling anti-patterns that waste time or create
  false confidence.
\end{itemize}

\section{The Core Model}

A blockchain threat model connects seven things:

\begin{lstlisting}
System:
  What are we reviewing?

Assets:
  What can be lost, corrupted, frozen, diluted, censored, or misrepresented?

Actors:
  Who can interact with the system, operate it, influence it, or attack it?

Entry points:
  Which transactions, messages, functions, instructions, proofs, admin actions,
  APIs, bots, and human processes can change state or influence decisions?

Assumptions:
  What must be true outside the local code for the security claim to hold?

Threats:
  What can go wrong if actors use entry points adversarially or assumptions fail?

Response:
  What will we build, test, monitor, document, or explicitly accept?
\end{lstlisting}

In compact form:

\begin{lstlisting}
Threat = actor + capability + entry point + vulnerable assumption + impact
\end{lstlisting}

Example:

\begin{lstlisting}
Actor:
  MEV searcher with flash liquidity

Capability:
  Can manipulate a low-liquidity AMM pool for one block

Entry point:
  borrow(collateral, amount)

Vulnerable assumption:
  Oracle spot price reflects fair collateral value

Impact:
  Attacker borrows against inflated collateral and leaves bad debt
\end{lstlisting}

That is a useful threat. It points to code, state, economic conditions,
tests, and mitigations.

Compare:

\begin{lstlisting}
Threat:
  Oracle attack.
\end{lstlisting}

That is a label. It is not yet a model.

\section{Turning a Protocol Description into a Security Model}

\subsection{Start With Scope}

Threat modeling begins by defining what is inside the review.

Scope may include:

\begin{itemize}
\tightlist
\item
  Smart contracts or programs.
\item
  Protocol docs and whitepaper.
\item
  Token contracts.
\item
  Governance contracts.
\item
  Upgrade mechanisms.
\item
  Deployed addresses.
\item
  Admin multisigs.
\item
  Frontend and wallet interactions.
\item
  Backend signers.
\item
  Keepers and bots.
\item
  Oracles.
\item
  Bridges and cross-chain messages.
\item
  Rollup or appchain assumptions.
\item
  Monitoring and incident response.
\item
  Human operational processes.
\end{itemize}

Scope should also state what is outside the review.

Bad:

\begin{lstlisting}
Review the protocol.
\end{lstlisting}

Better:

\begin{lstlisting}
Review the EVM vault, strategy adapters, share accounting, oracle integration,
upgrade path, and emergency controls. Exclude the frontend implementation and
the external lending market used by one strategy, but document both as trust
dependencies.
\end{lstlisting}

If the scope is unclear, the threat model will be unclear. If the threat
model is unclear, review effort will drift toward whatever is easiest to
inspect instead of whatever is most dangerous.

\subsection{Collect the Actual System, Not the Marketing System}

The system is not only the whitepaper.

Inputs to a real threat model include:

\begin{itemize}
\tightlist
\item
  Whitepaper.
\item
  Protocol docs.
\item
  Architecture diagrams.
\item
  Smart contract code.
\item
  Program code.
\item
  ABI/IDL/interface definitions.
\item
  Deployment addresses.
\item
  Proxy admin addresses.
\item
  Governance contracts.
\item
  Multisig owners and thresholds.
\item
  Timelock configurations.
\item
  Oracle configuration.
\item
  Token lists and accepted collateral.
\item
  Backend services.
\item
  Keeper and relayer design.
\item
  Frontend transaction construction.
\item
  Monitoring alerts.
\item
  Incident response plan.
\item
  Historical governance proposals.
\item
  Previous audit reports.
\item
  Known incidents and postmortems.
\end{itemize}

The docs describe the intended system. The deployment describes the
actual system. The code describes executable rules. Operations describe
who can change or rescue the system. Threat modeling needs all four.

If the docs say:

\begin{lstlisting}
Governance controls upgrades through a 48-hour timelock.
\end{lstlisting}

Verify:

\begin{lstlisting}
Which governance contract?
Which timelock?
Which proxy admin?
Which deployed implementation?
Who can cancel?
Who can execute?
Can anyone bypass the timelock?
Can emergency powers pause withdrawals?
\end{lstlisting}

Threat modeling starts polite. It does not stay naive.

\subsection{Extract Security Claims}

A protocol description usually contains security claims, even when it
does not call them that.

Examples:

\begin{lstlisting}
Depositors can always withdraw their fair share.
The oracle prevents price manipulation.
The bridge is trust-minimized.
Governance is decentralized.
The protocol is non-custodial.
The vault is delta-neutral.
Liquidations protect solvency.
The sequencer cannot steal funds.
The admin can only pause in emergencies.
The stablecoin is fully backed.
\end{lstlisting}

Convert each claim into reviewable form.

Weak:

\begin{lstlisting}
The protocol is non-custodial.
\end{lstlisting}

Reviewable:

\begin{lstlisting}
No privileged actor can transfer user funds, change withdrawal accounting,
upgrade to logic that transfers user funds, freeze withdrawals indefinitely,
or alter token mappings in a way that captures user claims without a documented
trust assumption and user exit window.
\end{lstlisting}

Weak:

\begin{lstlisting}
The bridge is trust-minimized.
\end{lstlisting}

Reviewable:

\begin{lstlisting}
Destination-chain minting requires a proof against finalized source-chain state,
source and destination domains are bound into the message, each message nonce is
consumed exactly once, and no multisig or upgrade path can mint unbacked assets
without an explicit trust assumption.
\end{lstlisting}

Every serious security claim should become either:

\begin{itemize}
\tightlist
\item
  An invariant.
\item
  A trust assumption.
\item
  A test.
\item
  A proof obligation.
\item
  A monitoring requirement.
\item
  A risk acceptance.
\end{itemize}

If a claim becomes none of those, it is probably marketing.

\subsection{Translate Features Into State Transitions}

Feature lists hide attack surfaces.

Example feature list:

\begin{lstlisting}
Users can deposit, withdraw, borrow, repay, and claim rewards.
Admins can configure markets.
Keepers handle liquidations.
Governance controls upgrades.
\end{lstlisting}

State-machine translation:

{\def\LTcaptype{none} % do not increment counter
\begin{longtable}[]{@{}
  >{\raggedright\arraybackslash}p{(\linewidth - 2\tabcolsep) * \real{0.5000}}
  >{\raggedright\arraybackslash}p{(\linewidth - 2\tabcolsep) * \real{0.5000}}@{}}
\toprule\noalign{}
\begin{minipage}[b]{\linewidth}\raggedright
Feature
\end{minipage} & \begin{minipage}[b]{\linewidth}\raggedright
State transitions
\end{minipage} \\
\midrule\noalign{}
\endhead
\bottomrule\noalign{}
\endlastfoot
Deposit & Transfer collateral in, mint shares or credit balance, update
indexes \\
Withdraw & Burn shares or debit balance, transfer collateral out, update
queues \\
Borrow & Read price, check collateral, mint or transfer debt asset,
update debt \\
Repay & Transfer debt asset in, reduce debt, update interest indexes \\
Claim rewards & Calculate accrued rewards, mark claimed, transfer reward
asset \\
Configure markets & Change collateral factors, caps, oracle, interest
model \\
Liquidate & Read price, check health, seize collateral, repay debt, pay
incentive \\
Upgrade & Change implementation or future transition rules \\
\end{longtable}
}

Now the reviewer can ask:

\begin{lstlisting}
Which transition touches assets?
Which transition reads external data?
Which transition can become insolvent?
Which transition can be called by anyone?
Which transition depends on ordering?
Which transition can be paused?
Which transition cannot be paused?
\end{lstlisting}

Threat modeling should turn nouns into transitions.

\subsection{Build a System Context Diagram}

You do not need a beautiful diagram. You need an accurate one.

A minimal blockchain system context diagram should show:

\begin{lstlisting}
Users / wallets
  -> frontend
  -> RPC provider
  -> chain
  -> protocol contracts/programs

Protocol contracts/programs
  -> tokens
  -> oracles
  -> external protocols
  -> governance
  -> admin multisig
  -> upgrade proxy
  -> keepers/relayers
  -> indexers/analytics
  -> monitoring/alerts
\end{lstlisting}

For cross-domain systems:

\begin{lstlisting}
source chain
  -> source bridge contract
  -> event/state proof
  -> relayer or validator set
  -> destination bridge contract
  -> wrapped token
\end{lstlisting}

For each arrow, ask:

\begin{lstlisting}
What crosses this boundary?
Who can modify it?
Who verifies it?
What happens if it is delayed, missing, stale, reordered, or false?
\end{lstlisting}


A simple attack tree turns a high-level threat into review paths:

\begin{lstlisting}
Goal: drain or disable the protocol
  -> abuse accounting transition
  -> manipulate oracle input
  -> compromise privileged authority
  -> exploit external call/control flow
  -> block withdrawals or exits
  -> exploit upgrade or governance path
\end{lstlisting}

Use the tree to decide where review time goes first. A threat model that does not change review priorities is only decoration.

The Threat Modeling Manifesto explicitly warns against looking for a
perfect representation and recommends multiple representations because
different views reveal different problems.\footnote{Threat Modeling
  Manifesto, https://www.threatmodelingmanifesto.org/.} For blockchain
protocols, useful views include:

\begin{itemize}
\tightlist
\item
  State-transition view.
\item
  Asset-flow view.
\item
  Authority graph.
\item
  Trust-boundary view.
\item
  Dependency graph.
\item
  Transaction-sequence view.
\item
  Upgrade/governance control graph.
\end{itemize}

No single diagram is sacred. The question is whether the representation
exposes risk.

\subsection{Threat Modeling Is Iterative}

OWASP emphasizes that threat modeling should be refined through the
lifecycle and updated after new features, incidents, or architecture
changes.\footnote{OWASP, ``Threat Modeling'',
  https://owasp.org/www-community/Threat\_Modeling.}

Blockchain systems especially require iteration because:

\begin{itemize}
\tightlist
\item
  Governance can change parameters.
\item
  New collateral can be added.
\item
  Oracles can be replaced.
\item
  Token behavior can differ across integrations.
\item
  Upgrade paths can change implementation.
\item
  TVL can grow and change attacker incentives.
\item
  Liquidity can move.
\item
  MEV conditions can change.
\item
  Bridge or rollup dependencies can change.
\item
  New chains and deployments can be added.
\end{itemize}

A threat model written before deployment may be obsolete after the first
governance proposal.

Threat model maintenance triggers:

\begin{itemize}
\tightlist
\item
  New asset listing.
\item
  New strategy.
\item
  New oracle.
\item
  New chain deployment.
\item
  Upgrade.
\item
  Governance change.
\item
  Admin role rotation.
\item
  Incident or near miss.
\item
  Large TVL increase.
\item
  Market regime change.
\item
  External dependency change.
\item
  New integration by another protocol.
\end{itemize}

If the threat model is not updated when the system changes, it becomes a
museum piece.

\section{Mapping Actors, Assets, Entry Points, and Privileged Roles}

\subsection{Actor Inventory}

Actors are entities that can interact with, operate, influence, or
depend on the system.

In blockchain systems, actors often include:

\begin{itemize}
\tightlist
\item
  End users.
\item
  Wallets.
\item
  Traders.
\item
  LPs.
\item
  Borrowers.
\item
  Liquidators.
\item
  Arbitrageurs.
\item
  MEV searchers.
\item
  Validators, miners, builders, and sequencers.
\item
  Governance token holders.
\item
  Delegates.
\item
  Admin multisigs.
\item
  Guardians.
\item
  Oracles and oracle operators.
\item
  Keepers.
\item
  Relayers.
\item
  Bridge validators.
\item
  Provers.
\item
  Frontend maintainers.
\item
  RPC providers.
\item
  Indexers.
\item
  Backend operators.
\item
  Custodians.
\item
  Exchanges.
\item
  Token issuers.
\item
  Auditors.
\item
  Attackers.
\item
  Insiders.
\end{itemize}

For each actor, record:

\begin{lstlisting}
Actor:
  Who or what is this?

Role:
  What are they supposed to do?

Authority:
  What transitions can they cause?

Capabilities:
  What resources, timing, information, or access do they have?

Incentives:
  What do they gain or lose?

Failure mode:
  What happens if they are malicious, compromised, unavailable, irrational,
  negligent, or mistaken?
\end{lstlisting}

This is intentionally broader than ``who can call functions.'' A
frontend maintainer may not call a privileged function, but can
construct malicious transactions. An RPC provider may not move funds,
but can distort a user's view. A keeper may not own assets, but can be
necessary for liveness.

\subsection{Asset Inventory}

Chapter 2 defined assets as enforceable claims over state. Threat
modeling requires a structured asset inventory.

Include:

\begin{itemize}
\tightlist
\item
  Native assets.
\item
  Tokens.
\item
  Shares.
\item
  LP positions.
\item
  Debt positions.
\item
  Collateral.
\item
  Pending withdrawals.
\item
  Rewards.
\item
  Bridge messages.
\item
  Governance power.
\item
  Upgrade authority.
\item
  Oracle authority.
\item
  Validator or sequencer keys.
\item
  Backend signing keys.
\item
  Domain names and frontend distribution.
\item
  Deployment secrets.
\item
  Monitoring credentials.
\item
  Off-chain reserves.
\item
  Legal claims.
\end{itemize}

For each asset:

\begin{lstlisting}
Asset:
  What has value?

Representation:
  Where is it recorded?

Normal transition:
  How should it move or change?

Privileged transition:
  Who can move, freeze, mint, burn, configure, or upgrade around it?

Invariant:
  What must remain true?

Impact:
  What happens if it is stolen, frozen, diluted, mispriced, duplicated,
  or made unrecoverable?
\end{lstlisting}

Do not restrict the asset inventory to user funds. Privileged keys are
often the most valuable assets in the system.

\subsection{Entry Point Inventory}

Entry points are ways state can be changed or security-relevant
decisions can be influenced.

On-chain entry points:

\begin{itemize}
\tightlist
\item
  Public functions.
\item
  Program instructions.
\item
  Module entry functions.
\item
  Governance proposal execution.
\item
  Upgrade functions.
\item
  Admin setters.
\item
  Oracle updates.
\item
  Bridge message handlers.
\item
  Permit or signature-based functions.
\item
  Callback hooks.
\item
  Token receiver functions.
\item
  Fallback and receive functions.
\item
  Batch and multicall paths.
\item
  Liquidation functions.
\item
  Emergency controls.
\end{itemize}

Off-chain entry points:

\begin{itemize}
\tightlist
\item
  Frontend transaction construction.
\item
  RPC request/response.
\item
  Indexer ingestion.
\item
  Backend APIs.
\item
  Keeper bots.
\item
  Relayer endpoints.
\item
  Signing services.
\item
  CI/CD pipelines.
\item
  Package publishing.
\item
  Governance forums.
\item
  Multisig signing workflows.
\item
  DNS and domain management.
\end{itemize}

For each entry point:

\begin{lstlisting}
Entry point:
  Function, instruction, API, workflow, or process.

Caller:
  Who can invoke it?

Authority:
  What authorization is required?

State read:
  What state does it depend on?

State written:
  What state can it change?

External assumptions:
  What oracle, token, RPC, relayer, signer, or human process does it trust?

Failure behavior:
  What happens on revert, partial failure, timeout, or non-response?

Abuse:
  What can an adversarial caller do with valid inputs?
\end{lstlisting}

This inventory becomes the review map. It tells you where to spend
attention.

\subsection{Privileged Role Inventory}

Privileged roles are authorities that can change rules, parameters,
dependencies, or emergency state.

Examples:

\begin{itemize}
\tightlist
\item
  Owner.
\item
  Proxy admin.
\item
  Upgrade executor.
\item
  Timelock proposer.
\item
  Timelock canceller.
\item
  Governance executor.
\item
  Pauser.
\item
  Guardian.
\item
  Minter.
\item
  Burner.
\item
  Blacklist manager.
\item
  Oracle admin.
\item
  Risk parameter admin.
\item
  Strategy manager.
\item
  Fee recipient setter.
\item
  Bridge validator-set admin.
\item
  Sequencer admin.
\item
  Backend signer.
\item
  Deployment key.
\end{itemize}

For each privileged role:

\begin{lstlisting}
Role:
  Name and contract/program/module.

Holder:
  EOA, multisig, contract, governance, DAO, backend service, or operator.

Capability:
  Exact transitions authorized.

Delay:
  Timelock, voting delay, challenge window, or no delay.

Bypass:
  Emergency path, guardian override, admin override, or upgrade path.

Rotation:
  How can it be changed?

Loss:
  What happens if lost?

Compromise:
  What happens if compromised?

Observability:
  Can users see pending action before harm?
\end{lstlisting}

Role names lie by omission. A ``guardian'' may only pause deposits, or
it may pause withdrawals and trap funds. A ``risk admin'' may only lower
caps, or it may raise collateral factors to unsafe values. A
``governor'' may only adjust fees, or it may call arbitrary targets.

Write the verbs.

\subsection{Data and Control Flows}

Traditional threat modeling often uses data-flow diagrams and trust
boundaries. Microsoft STRIDE guidance frames threats around questions
such as whether authentication data can be changed, whether data can be
read, or whether access can be denied, and categorizes threats as
spoofing, tampering, repudiation, information disclosure, denial of
service, and elevation of privilege.\footnote{Microsoft, ``Microsoft
  Threat Modeling Tool threats'',
  https://learn.microsoft.com/en-us/azure/security/develop/threat-modeling-tool-threats.}

Blockchain systems need both data-flow and control-flow views.

Data flows:

\begin{itemize}
\tightlist
\item
  Oracle price -\textgreater{} lending market.
\item
  Event log -\textgreater{} bridge relayer -\textgreater{} destination
  chain.
\item
  RPC response -\textgreater{} frontend -\textgreater{} wallet.
\item
  Governance proposal JSON -\textgreater{} forum -\textgreater{} UI
  -\textgreater{} voters.
\item
  Indexer state -\textgreater{} dashboard -\textgreater{} user decision.
\item
  Transaction simulation -\textgreater{} wallet display.
\item
  Keeper observation -\textgreater{} liquidation transaction.
\end{itemize}

Control flows:

\begin{itemize}
\tightlist
\item
  User transaction -\textgreater{} contract call -\textgreater{}
  external token hook -\textgreater{} callback.
\item
  Governance vote -\textgreater{} timelock -\textgreater{} proxy admin
  -\textgreater{} implementation change.
\item
  Admin setter -\textgreater{} oracle address -\textgreater{}
  price-dependent borrow path.
\item
  Sequencer ordering -\textgreater{} liquidation execution
  -\textgreater{} solvency.
\item
  Multisig approval -\textgreater{} upgrade -\textgreater{} storage
  layout migration.
\item
  Relayer submission -\textgreater{} proof verification -\textgreater{}
  mint.
\end{itemize}

For each flow:

\begin{lstlisting}
What crosses the boundary?
Who can alter it?
Who verifies it?
What state transition depends on it?
What if it is stale, false, delayed, missing, or reordered?
\end{lstlisting}

\subsection{The Minimum Map}

Before reviewing a serious protocol, you should be able to fill this
table.

{\def\LTcaptype{none} % do not increment counter
\begin{longtable}[]{@{}
  >{\raggedright\arraybackslash}p{(\linewidth - 2\tabcolsep) * \real{0.5000}}
  >{\raggedright\arraybackslash}p{(\linewidth - 2\tabcolsep) * \real{0.5000}}@{}}
\toprule\noalign{}
\begin{minipage}[b]{\linewidth}\raggedright
Category
\end{minipage} & \begin{minipage}[b]{\linewidth}\raggedright
Minimum questions
\end{minipage} \\
\midrule\noalign{}
\endhead
\bottomrule\noalign{}
\endlastfoot
Assets & What can be stolen, frozen, diluted, mispriced, or made
unrecoverable? \\
Actors & Who can interact, operate, influence, or attack? \\
Entry points & What can change state or influence a state-changing
decision? \\
Privileged roles & Who can change parameters, dependencies, code, or
emergency state? \\
Trust boundaries & What does the system rely on but not verify
locally? \\
External dependencies & Which tokens, oracles, chains, bridges,
relayers, RPCs, indexers, and humans matter? \\
Invariants & What must remain true across transitions? \\
Liveness paths & What required transitions must remain possible? \\
Monitoring & What must be observed to detect assumption failure? \\
\end{longtable}
}

If a team cannot answer these questions, the team is not ready for a
deep audit. It is still discovering the system.

\section{Separating Users, Attackers, Rational Participants, and Trusted Operators}

\subsection{Do Not Use One Generic Attacker}

``The attacker'' is often too broad to be useful.

Different adversaries have different powers:

\begin{itemize}
\tightlist
\item
  A normal user can call public entry points.
\item
  A bot can call them faster and more often.
\item
  A searcher can simulate pending transactions and pay for ordering.
\item
  A validator or sequencer can influence inclusion and ordering.
\item
  A governance whale can pass proposals.
\item
  An admin can call privileged functions.
\item
  A compromised frontend can construct malicious transactions.
\item
  A malicious token can behave unexpectedly during integration.
\item
  A bridge validator set can sign false messages.
\item
  A custodian can delay or block withdrawals.
\item
  An insider can access keys, deployments, secrets, or configs.
\item
  A state-level actor may have long time horizons, supply-chain reach,
  and legal or coercive power.
\end{itemize}

NIST SP 800-30 recommends identifying threat sources and characterizing
adversarial sources by factors such as capability, intent, and
targeting.\footnote{NIST, ``SP 800-30 Revision 1: Guide for Conducting
  Risk Assessments'', threat-source identification and characterization,
  https://nvlpubs.nist.gov/nistpubs/Legacy/SP/nistspecialpublication800-30r1.pdf.}

For blockchain systems, also add:

\begin{itemize}
\tightlist
\item
  Capital.
\item
  Liquidity.
\item
  Ordering access.
\item
  Governance power.
\item
  Cross-domain reach.
\item
  Operational access.
\item
  Ability to wait.
\item
  Ability to take reputational or legal risk.
\end{itemize}

Threat modeling improves when each actor has a realistic capability set.

\subsection{Users Can Be Adversarial Without Being ``Hackers''}

In permissionless systems, any user may call public functions with valid
inputs. That does not make them a hacker. It makes them a participant.

Examples:

\begin{itemize}
\tightlist
\item
  A trader submits a sandwich.
\item
  A liquidator triggers liquidation at the worst moment for a borrower.
\item
  A borrower maximizes leverage.
\item
  A depositor withdraws just before losses are realized.
\item
  A governance voter sells votes.
\item
  A user creates thousands of addresses.
\item
  A keeper ignores unprofitable upkeep.
\end{itemize}

These may be rational behaviors inside the rules.

Security model:

\begin{lstlisting}
Users are not trusted to follow product intent.
They are trusted only to satisfy validity rules and pay costs.
\end{lstlisting}

This is why business logic must be robust against adversarial but valid
use.

\subsection{Rational Participants Are Not Byzantine Actors}

A rational actor pursues payoff. A Byzantine actor may behave
arbitrarily, including actions that appear irrational, malicious,
inconsistent, or self-harming.

This distinction matters.

Rational searcher:

\begin{lstlisting}
Will reorder or backrun if profit exceeds cost.
Probably will not burn money without benefit.
\end{lstlisting}

Byzantine validator:

\begin{lstlisting}
May censor, equivocate, or violate protocol rules if capable,
even if immediate profit is unclear.
\end{lstlisting}

Compromised key:

\begin{lstlisting}
Can sign any message the key is authorized to sign.
The original owner's incentives no longer matter.
\end{lstlisting}

Negligent operator:

\begin{lstlisting}
May fail to update, rotate, monitor, or respond.
The failure may not be malicious.
\end{lstlisting}

The mitigation differs:

\begin{itemize}
\tightlist
\item
  Rational attacks require economic bounds.
\item
  Byzantine attacks require validation, quorum assumptions, proofs, or
  slashing.
\item
  Compromise requires key management and blast-radius reduction.
\item
  Negligence requires process controls, automation, monitoring, and
  recovery paths.
\end{itemize}

Do not prescribe one defense for different fault models.

\subsection{Trusted Operators Are Threat Sources Too}

Trusted operators are actors the system relies on:

\begin{itemize}
\tightlist
\item
  Admin multisig signers.
\item
  Governance delegates.
\item
  Oracle operators.
\item
  Bridge validators.
\item
  Sequencer operators.
\item
  Relayers.
\item
  Keepers.
\item
  Backend signers.
\item
  Frontend maintainers.
\item
  DevOps engineers.
\item
  Custodians.
\item
  Token issuers.
\end{itemize}

They may fail in several ways:

\begin{itemize}
\tightlist
\item
  Malicious action.
\item
  Key compromise.
\item
  Downtime.
\item
  Misconfiguration.
\item
  Slow response.
\item
  Collusion.
\item
  Legal compulsion.
\item
  Insider abuse.
\item
  Incorrect deployment.
\item
  Poor monitoring.
\end{itemize}

Treating trusted operators as outside the threat model is a classic
mistake.

Better:

\begin{lstlisting}
The oracle committee is trusted to report prices within acceptable deviation
and availability bounds. The protocol fails safely if one publisher is wrong,
but may become unsafe if a quorum signs a bad price or if updates stop for more
than 30 minutes during high volatility.
\end{lstlisting}

This names both safety and liveness assumptions.

\subsection{Insiders Have Different Access}

Insiders may have access that outsiders do not:

\begin{itemize}
\tightlist
\item
  Unreleased code.
\item
  Deployment keys.
\item
  Admin interfaces.
\item
  CI/CD secrets.
\item
  Monitoring credentials.
\item
  Backend databases.
\item
  Signing services.
\item
  Governance proposal drafts.
\item
  Oracle update infrastructure.
\item
  Private relayer endpoints.
\item
  Multisig coordination channels.
\item
  Incident response plans.
\end{itemize}

Insider threat modeling asks:

\begin{lstlisting}
What can one insider do?
What requires two insiders?
What is logged?
What is reviewed?
What can be done without on-chain visibility?
What can be done before users can react?
\end{lstlisting}

This is not paranoia. It is basic authority mapping.

\subsection{Actor Table Template}

Use this table.

{\def\LTcaptype{none} % do not increment counter
\begin{longtable}[]{@{}
  >{\raggedright\arraybackslash}p{(\linewidth - 10\tabcolsep) * \real{0.1667}}
  >{\raggedright\arraybackslash}p{(\linewidth - 10\tabcolsep) * \real{0.1667}}
  >{\raggedright\arraybackslash}p{(\linewidth - 10\tabcolsep) * \real{0.1667}}
  >{\raggedright\arraybackslash}p{(\linewidth - 10\tabcolsep) * \real{0.1667}}
  >{\raggedright\arraybackslash}p{(\linewidth - 10\tabcolsep) * \real{0.1667}}
  >{\raggedright\arraybackslash}p{(\linewidth - 10\tabcolsep) * \real{0.1667}}@{}}
\toprule\noalign{}
\begin{minipage}[b]{\linewidth}\raggedright
Actor
\end{minipage} & \begin{minipage}[b]{\linewidth}\raggedright
Normal role
\end{minipage} & \begin{minipage}[b]{\linewidth}\raggedright
Capabilities
\end{minipage} & \begin{minipage}[b]{\linewidth}\raggedright
Incentives
\end{minipage} & \begin{minipage}[b]{\linewidth}\raggedright
Failure mode
\end{minipage} & \begin{minipage}[b]{\linewidth}\raggedright
Controls
\end{minipage} \\
\midrule\noalign{}
\endhead
\bottomrule\noalign{}
\endlastfoot
User & Deposit/withdraw/borrow & Public calls, signatures &
Utility/profit & Adversarial sequencing, phishing & Input checks,
slippage, limits \\
Searcher & Extract MEV & Simulation, bundles, priority fees & Profit &
Ordering attack & MEV-aware design, private flow assumptions \\
Admin multisig & Emergency/upgrade & Privileged calls & Protocol
maintenance & Collusion, compromise, wrong signing & Timelock,
threshold, monitoring \\
Oracle & Price data & Signed updates & Correct reporting &
Stale/manipulated price & Deviation checks, circuit breakers \\
Keeper & Maintenance & Calls upkeep functions & Fees/rewards & No
liveness if unprofitable & Incentives, fallback, permissionless calls \\
\end{longtable}
}

Fill it with actual actors. Generic roles are placeholders, not
conclusions.

\section{Identifying Trust Assumptions and External Dependencies}

\subsection{Assumptions Are Security Debt}

Every system has assumptions. The problem is not having assumptions. The
problem is hiding them.

Examples:

\begin{lstlisting}
The oracle price is fresh.
The bridge source event is final.
The token does not reenter.
The admin will not upgrade maliciously.
The sequencer will not censor exits indefinitely.
The relayer will submit messages.
The keeper reward is enough.
The frontend points to the correct contracts.
The RPC provider reports current chain state.
The stablecoin issuer honors redemption.
\end{lstlisting}

Each assumption is a place where the system is not self-contained.

Write assumptions in a falsifiable form:

\begin{lstlisting}
The protocol assumes Chainlink ETH/USD updates at least every 1 hour
and that the returned answer is not stale, negative, or outside configured
deviation bounds.
\end{lstlisting}

Not:

\begin{lstlisting}
We use a secure oracle.
\end{lstlisting}

\subsection{Dependency Inventory}

Common blockchain dependencies:

{\def\LTcaptype{none} % do not increment counter
\begin{longtable}[]{@{}
  >{\raggedright\arraybackslash}p{(\linewidth - 4\tabcolsep) * \real{0.3333}}
  >{\raggedright\arraybackslash}p{(\linewidth - 4\tabcolsep) * \real{0.3333}}
  >{\raggedright\arraybackslash}p{(\linewidth - 4\tabcolsep) * \real{0.3333}}@{}}
\toprule\noalign{}
\begin{minipage}[b]{\linewidth}\raggedright
Dependency
\end{minipage} & \begin{minipage}[b]{\linewidth}\raggedright
Common assumption
\end{minipage} & \begin{minipage}[b]{\linewidth}\raggedright
Failure
\end{minipage} \\
\midrule\noalign{}
\endhead
\bottomrule\noalign{}
\endlastfoot
Base chain & Finality, liveness, valid execution & Reorg, halt,
censorship \\
Token contract & Standard behavior & Fee-on-transfer, blacklist, rebase,
hook \\
Oracle & Fresh and manipulation-resistant data & Bad price, stale price,
missing update \\
AMM/liquidity & Price reflects executable market & Low-liquidity
manipulation \\
Governance & Non-captured decision process & Malicious proposal,
bribery, low quorum \\
Admin multisig & Honest and available signers & Collusion, compromise,
lost keys \\
Bridge & Correct message verification & Forged messages, replay,
finality mismatch \\
Sequencer & Inclusion/liveness assumptions & Censorship, downtime,
ordering extraction \\
RPC & Correct view of chain & Incorrect responses, censorship, lag \\
Indexer & Correct interpretation & Desync, missed reorg, event
misread \\
Frontend & Correct transaction construction & Malicious approvals, wrong
addresses \\
Backend signer & Correct policy enforcement & Unauthorized
withdrawals \\
Keeper & Timely maintenance & Stuck positions, bad debt \\
CI/CD & Integrity of deployment & Malicious release, wrong address \\
Custodian & Asset safekeeping & Insolvency, freeze, withdrawal halt \\
\end{longtable}
}

Each row can become a threat scenario.

\subsection{Safety vs Liveness Assumptions}

Separate assumptions that preserve safety from assumptions that preserve
liveness.

Safety assumptions:

\begin{itemize}
\tightlist
\item
  Oracle price is not manipulable enough to cause undercollateralized
  borrowing.
\item
  Bridge validator quorum does not sign false messages.
\item
  Upgrade admin cannot bypass timelock.
\item
  Token mint authority cannot create unbacked supply.
\item
  Contract verifier rejects malformed proofs.
\end{itemize}

Liveness assumptions:

\begin{itemize}
\tightlist
\item
  At least one keeper will liquidate unhealthy positions.
\item
  Relayers will submit bridge messages.
\item
  Users can withdraw during an exit window.
\item
  Sequencer downtime fallback works.
\item
  Governance can execute emergency fixes.
\item
  Proof generation service remains available.
\end{itemize}

Some assumptions affect both:

\begin{itemize}
\tightlist
\item
  Sequencer censorship can block exits and affect liquidation ordering.
\item
  Oracle downtime can freeze borrowing and liquidations.
\item
  Admin key loss can prevent emergency response.
\item
  Bridge relayer compromise can become safety risk if relayers are
  trusted for message truth.
\end{itemize}

For each assumption, write:

\begin{lstlisting}
Type:
  Safety, liveness, or both.

Failure impact:
  What bad state or stuck state results?

Detection:
  How do we know it failed?

Response:
  What can the system do?
\end{lstlisting}

\subsection{External Contracts Are Dependencies, Not Background}

If your protocol calls another contract, integrates a token, deposits
into a strategy, reads an oracle, or relies on a registry, that external
system is inside the threat model as a dependency.

External contract risks:

\begin{itemize}
\tightlist
\item
  Upgradeable dependency.
\item
  Reentrancy or callbacks.
\item
  Nonstandard token behavior.
\item
  Incorrect decimals.
\item
  Pausable or blacklistable token.
\item
  Fee-on-transfer or rebasing token.
\item
  Governance-controlled parameter changes.
\item
  Oracle update delays.
\item
  Liquidity migration.
\item
  Contract selfdestruct or deprecated behavior.
\item
  Address substitution.
\end{itemize}

Dependency review questions:

\begin{lstlisting}
Can the dependency change behavior?
Who controls it?
Can it call back?
Can it return false, no data, malformed data, or stale data?
Can it pause or freeze?
Can it become illiquid?
Can it censor or reject this protocol?
Can it break accounting assumptions?
\end{lstlisting}

Do not integrate an external contract as if it were a math function.

\subsection{Human Dependencies Are Real Dependencies}

Human processes are often the hidden critical path:

\begin{itemize}
\tightlist
\item
  Someone notices an oracle failure.
\item
  Someone signs a pause transaction.
\item
  Someone rotates a compromised key.
\item
  Someone updates the frontend.
\item
  Someone submits the governance proposal.
\item
  Someone pays relayer gas.
\item
  Someone checks deployment addresses.
\item
  Someone verifies a multisig transaction.
\end{itemize}

Threat model them.

Questions:

\begin{lstlisting}
Who is responsible?
What signal triggers action?
How fast must action happen?
What if the person is unavailable?
What if signers disagree?
What if the runbook is wrong?
What if the attacker acts faster?
What if the emergency control itself is compromised?
\end{lstlisting}

If the system's safety depends on a human responding in five minutes at
3 AM, write that assumption plainly.

\subsection{Assumption Register}

Use an assumption register.

{\def\LTcaptype{none} % do not increment counter
\begin{longtable}[]{@{}
  >{\raggedright\arraybackslash}p{(\linewidth - 12\tabcolsep) * \real{0.1429}}
  >{\raggedright\arraybackslash}p{(\linewidth - 12\tabcolsep) * \real{0.1429}}
  >{\raggedright\arraybackslash}p{(\linewidth - 12\tabcolsep) * \real{0.1429}}
  >{\raggedright\arraybackslash}p{(\linewidth - 12\tabcolsep) * \real{0.1429}}
  >{\raggedright\arraybackslash}p{(\linewidth - 12\tabcolsep) * \real{0.1429}}
  >{\raggedright\arraybackslash}p{(\linewidth - 12\tabcolsep) * \real{0.1429}}
  >{\raggedright\arraybackslash}p{(\linewidth - 12\tabcolsep) * \real{0.1429}}@{}}
\toprule\noalign{}
\begin{minipage}[b]{\linewidth}\raggedright
ID
\end{minipage} & \begin{minipage}[b]{\linewidth}\raggedright
Assumption
\end{minipage} & \begin{minipage}[b]{\linewidth}\raggedright
Type
\end{minipage} & \begin{minipage}[b]{\linewidth}\raggedright
Dependency
\end{minipage} & \begin{minipage}[b]{\linewidth}\raggedright
Failure impact
\end{minipage} & \begin{minipage}[b]{\linewidth}\raggedright
Detection
\end{minipage} & \begin{minipage}[b]{\linewidth}\raggedright
Mitigation
\end{minipage} \\
\midrule\noalign{}
\endhead
\bottomrule\noalign{}
\endlastfoot
A1 & Oracle updates within 30 minutes & Safety/liveness & Oracle network
& Bad debt or market freeze & Price age alert & Staleness checks,
pause \\
A2 & 3-of-5 admin multisig not compromised & Safety & Signers &
Malicious upgrade & Timelock monitor & Higher threshold, delay \\
A3 & Relayer submits messages within 1 hour & Liveness & Relayer network
& Stuck withdrawals & Queue age alert & Permissionless relay \\
A4 & Token has no transfer fee & Safety & Collateral token & Accounting
drift & Balance delta check & Measure actual received \\
\end{longtable}
}

An assumption register is valuable because it gives reviewers a target.
Each assumption can be challenged, tested, monitored, or accepted
explicitly.

\section{Mapping Attack Surfaces Across Chain, Contract, Infrastructure, and Human Layers}

\subsection{Blockchain Attack Surface Is Layered}

Blockchain security fails across layers:

\begin{lstlisting}
Chain layer:
  Consensus, finality, reorgs, censorship, block production.

Transaction/order layer:
  Mempools, MEV, fees, ordering, inclusion, replacement.

Execution layer:
  VM semantics, gas/compute limits, account model, callbacks.

Contract/protocol layer:
  Authorization, accounting, external calls, upgrades, invariants.

Economic layer:
  Liquidity, oracle manipulation, incentives, governance capture.

Cross-domain layer:
  Bridges, rollups, relayers, proof systems, finality mismatch.

Infrastructure layer:
  RPCs, indexers, frontends, backends, CI/CD, dependencies.

Human layer:
  Multisig signing, incident response, custody, social engineering.
\end{lstlisting}

Many real exploits cross layers.

Example:

\begin{lstlisting}
Low-liquidity AMM price
  -> oracle manipulation
  -> lending borrow
  -> bad debt
  -> governance emergency pause
  -> users unable to withdraw
\end{lstlisting}

If your threat model only covers Solidity code, you miss the exploit
path.

\subsection{STRIDE as a Prompt, Not a Cage}

STRIDE is useful as a prompt:

\begin{itemize}
\tightlist
\item
  Spoofing.
\item
  Tampering.
\item
  Repudiation.
\item
  Information disclosure.
\item
  Denial of service.
\item
  Elevation of privilege.
\end{itemize}

Microsoft describes STRIDE as a way to categorize threats and ask
pointed questions about authentication, data modification, data
exposure, service denial, and privilege escalation.\footnote{Microsoft,
  ``Microsoft Threat Modeling Tool threats'',
  https://learn.microsoft.com/en-us/azure/security/develop/threat-modeling-tool-threats.}

Blockchain adaptation:

{\def\LTcaptype{none} % do not increment counter
\begin{longtable}[]{@{}
  >{\raggedright\arraybackslash}p{(\linewidth - 2\tabcolsep) * \real{0.5000}}
  >{\raggedright\arraybackslash}p{(\linewidth - 2\tabcolsep) * \real{0.5000}}@{}}
\toprule\noalign{}
\begin{minipage}[b]{\linewidth}\raggedright
STRIDE category
\end{minipage} & \begin{minipage}[b]{\linewidth}\raggedright
Blockchain examples
\end{minipage} \\
\midrule\noalign{}
\endhead
\bottomrule\noalign{}
\endlastfoot
Spoofing & Forged signer, fake token, fake frontend, fake bridge
message, fake oracle \\
Tampering & Storage corruption, malicious upgrade, manipulated price,
altered calldata \\
Repudiation & Unclear signer intent, missing audit logs, unverifiable
admin process \\
Information disclosure & Leaked private order flow, exposed xpub, leaked
signing policy, metadata deanonymization \\
Denial of service & Stuck withdrawals, gas griefing, account locking,
sequencer censorship, proof too large \\
Elevation of privilege & Missing access control, role-admin abuse,
governance capture, proxy admin takeover \\
\end{longtable}
}

STRIDE should not replace blockchain-specific reasoning. It should help
you ask what can go wrong from several angles.

\subsection{Contract and Program Attack Surface}

Contract/program surfaces include:

\begin{itemize}
\tightlist
\item
  Public functions and instructions.
\item
  External calls.
\item
  Delegatecall or equivalent dynamic dispatch.
\item
  Token transfers.
\item
  Callbacks and hooks.
\item
  Signature verification.
\item
  Access control.
\item
  Upgradeability.
\item
  Initialization.
\item
  Storage layout.
\item
  Math and rounding.
\item
  Oracle reads.
\item
  Event emission.
\item
  Batch/multicall.
\item
  Fallback logic.
\item
  Account constraints.
\item
  Resource limits.
\end{itemize}

OWASP's Smart Contract Top 10 includes categories such as access
control, business logic, oracle manipulation, flash-loan-facilitated
attacks, input validation, unchecked external calls, arithmetic errors,
reentrancy, and upgradeability risk.\footnote{OWASP, ``Smart Contract
  Top 10'', https://owasp.org/www-project-smart-contract-top-10/.}

Use such lists as prompts, not as a finish line. A protocol can avoid
every named category and still fail economically because its invariant
is wrong.

\subsection{Chain and Ordering Attack Surface}

Chain-layer and ordering surfaces include:

\begin{itemize}
\tightlist
\item
  Reorgs.
\item
  Delayed finality.
\item
  Chain halts.
\item
  Censorship.
\item
  Mempool visibility.
\item
  Transaction replacement.
\item
  Priority fees.
\item
  Builder/searcher behavior.
\item
  Sequencer downtime.
\item
  Forced inclusion.
\item
  Cross-domain settlement delay.
\end{itemize}

Threat questions:

\begin{lstlisting}
Can a profitable actor change ordering?
Can a transaction be sandwiched?
Can a liquidation be delayed?
Can a bridge event be reorged out?
Can exits be censored?
Can a user rely on state before settlement?
Can a governance proposal be front-run or blocked?
\end{lstlisting}

If the protocol depends on timely or private execution, ordering is part
of the threat model.

\subsection{Economic Attack Surface}

Economic security is not optional in blockchain systems because assets
and incentives are native to the execution environment.

Economic surfaces include:

\begin{itemize}
\tightlist
\item
  Liquidity depth.
\item
  Oracle reference markets.
\item
  Collateral factors.
\item
  Liquidation incentives.
\item
  Auction design.
\item
  Reward emissions.
\item
  Governance token distribution.
\item
  Vote borrowing and bribery.
\item
  Peg mechanisms.
\item
  Withdrawal queues.
\item
  Restaking correlation.
\item
  Insurance funds.
\item
  Bad debt handling.
\end{itemize}

Threat questions:

\begin{lstlisting}
How much capital does the attack require?
Can flash liquidity reduce that capital requirement?
Can the attacker manipulate a price only briefly?
Can the attacker profit atomically?
Who absorbs losses?
Can a rational actor grief for indirect profit?
Can incentives fail during market stress?
\end{lstlisting}

Economic attack surface cannot be found by scanning code alone.

\subsection{Infrastructure Attack Surface}

Infrastructure surfaces include:

\begin{itemize}
\tightlist
\item
  Frontend hosting.
\item
  DNS.
\item
  Wallet connectors.
\item
  RPC providers.
\item
  Indexers.
\item
  Backend APIs.
\item
  Signing services.
\item
  CI/CD.
\item
  Package dependencies.
\item
  Container images.
\item
  Monitoring tools.
\item
  Admin dashboards.
\item
  Cloud accounts.
\item
  Secrets management.
\end{itemize}

OWASP's Web3 attack-vector material explicitly goes beyond smart
contracts and includes risks such as multisig hijacking, supply-chain
attacks, private key compromise, UI/UX spoofing and approval phishing,
centralized infrastructure breaches, DNS and routing hijacking, wallet
compromise, insider threats, and malicious open-source
contributions.\footnote{OWASP Smart Contract Security, ``Alternate Top
  15 - Web3 Attack Vectors Beyond Smart Contracts'',
  https://scs.owasp.org/sctop10/Web3-Attack-Vectors-Top15/.}

This matters because many ``protocol'' losses happen outside protocol
code.

Threat questions:

\begin{lstlisting}
Can a compromised frontend construct malicious approvals?
Can DNS hijack redirect users?
Can an indexer desync make users take unsafe actions?
Can backend signers authorize withdrawals?
Can CI/CD deploy wrong contracts?
Can a package update exfiltrate keys?
Can RPC censorship hide pending danger?
\end{lstlisting}

If an off-chain component can cause users or operators to authorize
harmful transitions, it is in scope.

\subsection{Human and Governance Attack Surface}

Human surfaces include:

\begin{itemize}
\tightlist
\item
  Multisig review.
\item
  Key backups.
\item
  Deployment ceremonies.
\item
  Governance discussion.
\item
  Vote delegation.
\item
  Emergency response.
\item
  Legal and compliance controls.
\item
  Operator incentives.
\item
  Public communications.
\item
  Support channels.
\item
  Social engineering.
\end{itemize}

Governance surfaces include:

\begin{itemize}
\tightlist
\item
  Proposal thresholds.
\item
  Quorum.
\item
  Voting power distribution.
\item
  Delegation.
\item
  Timelock.
\item
  Execution payloads.
\item
  Cancellation rights.
\item
  Emergency veto.
\item
  Treasury control.
\item
  Governance self-upgrade.
\end{itemize}

Threat questions:

\begin{lstlisting}
Can governance call arbitrary targets?
Can voting power be borrowed?
Can delegates be bribed?
Can a proposal hide malicious calldata?
Can signers verify what they approve?
Can an attacker exploit public incident response?
Can a social-engineered employee deploy malicious code?
\end{lstlisting}

The human layer often holds the keys to the machine.

\section{Producing a Threat Model That Guides Review Instead of Decorating It}

\subsection{The Threat Model Must Produce Work}

A useful threat model produces review tasks.

For each important threat, the model should say:

\begin{lstlisting}
Threat:
  What can go wrong?

Priority:
  Why does it matter more or less than others?

Review target:
  Which code, config, process, or dependency must be inspected?

Invariant:
  What property should be written or tested?

Mitigation:
  What design control prevents or limits the threat?

Residual risk:
  What remains after mitigation?

Monitoring:
  How would we detect failure?
\end{lstlisting}

If a threat model does not affect review priorities, it is probably
decoration.

\subsection{Prioritize by Impact, Feasibility, and Uncertainty}

NIST SP 800-30 frames risk around threat events, threat sources,
vulnerabilities or predisposing conditions, likelihood, and
impact.\footnote{NIST, ``SP 800-30 Revision 1: Guide for Conducting Risk
  Assessments'', likelihood determination,
  https://nvlpubs.nist.gov/nistpubs/Legacy/SP/nistspecialpublication800-30r1.pdf.}
It also notes that impact determinations should make assumptions,
sources, methods, and rationale explicit.\footnote{NIST, ``SP 800-30
  Revision 1: Guide for Conducting Risk Assessments'', impact
  determination,
  https://nvlpubs.nist.gov/nistpubs/Legacy/SP/nistspecialpublication800-30r1.pdf.}

For blockchain security review, use practical factors:

\begin{itemize}
\tightlist
\item
  Funds at risk.
\item
  Privilege impact.
\item
  Solvency impact.
\item
  Governance impact.
\item
  Liveness impact.
\item
  Cross-domain contagion.
\item
  Attacker capital required.
\item
  Exploit reliability.
\item
  Preconditions.
\item
  Time sensitivity.
\item
  Detectability.
\item
  Recovery difficulty.
\item
  Uncertainty.
\end{itemize}

Uncertainty deserves emphasis. A high-impact area with unclear
assumptions should get review attention even before a known bug exists.

Risk shorthand:

\begin{lstlisting}
High priority:
  High impact + plausible path + unclear or weak controls.

Medium priority:
  Meaningful impact + constrained path or strong controls.

Low priority:
  Low impact, implausible path, or well-constrained residual risk.

Research priority:
  Impact potentially high but model uncertain.
\end{lstlisting}

Do not pretend you can reduce every blockchain threat to a neat numeric
score. Use scoring to communicate, not to replace judgment.

\subsection{Turn Threats Into Invariants and Tests}

Threat:

\begin{lstlisting}
Attacker manipulates vault share price by donating assets before a victim deposit.
\end{lstlisting}

Invariant:

\begin{lstlisting}
Depositing assets should mint shares proportional to economic contribution,
and no direct asset transfer should let an attacker force later deposits to mint
zero or unfairly low shares.
\end{lstlisting}

Tests:

\begin{lstlisting}
first depositor fuzzing
direct donation before deposit
small deposit rounding
large asset balance / tiny share supply
differential comparison with ideal math model
\end{lstlisting}

Threat:

\begin{lstlisting}
Admin upgrades vault to malicious implementation during timelock.
\end{lstlisting}

Invariant or security property:

\begin{lstlisting}
Users can withdraw before a delayed upgrade executes, unless withdrawals are
paused only under documented emergency rules.
\end{lstlisting}

Tests/review:

\begin{lstlisting}
simulate upgrade queue and execution
verify timelock delay
verify pause cannot block exit indefinitely without escalation
verify monitoring can detect queued implementation
\end{lstlisting}

The threat model should generate the testing plan.

\subsection{Decide Mitigation, Acceptance, Transfer, or Elimination}

Not every risk is fixed the same way.

Possible responses:

\begin{itemize}
\tightlist
\item
  Eliminate the feature.
\item
  Redesign the transition.
\item
  Add validation.
\item
  Add invariant checks.
\item
  Add rate limits.
\item
  Add timelock.
\item
  Add circuit breaker.
\item
  Add monitoring.
\item
  Restrict supported assets.
\item
  Add independent oracle.
\item
  Add fallback path.
\item
  Increase multisig threshold.
\item
  Improve signing workflow.
\item
  Add tests or formal specifications.
\item
  Document and accept residual risk.
\item
  Transfer risk through insurance or legal process.
\end{itemize}

OWASP's threat-modeling guidance includes identifying countermeasures or
managing risk after identifying what can go wrong.\footnote{OWASP,
  ``Threat Modeling'', https://owasp.org/www-community/Threat\_Modeling.}

For each material threat, write the decision:

\begin{lstlisting}
Mitigate:
  Add staleness and deviation checks.

Accept:
  Governance can change collateral factor after timelock; users accept governance risk.

Eliminate:
  Do not support fee-on-transfer collateral.

Monitor:
  Alert if oracle age > 30 minutes or market price deviates > 2 percent.
\end{lstlisting}

Unwritten acceptance is not risk management. It is amnesia.

\subsection{Threat Model Deliverable Template}

A useful threat model can be concise if it is structured.

\begin{lstlisting}
1. Scope
  Included systems:
  Excluded systems:
  Review goals:

2. System summary
  State machines:
  Major flows:
  Deployment addresses:

3. Assets
  User assets:
  Protocol assets:
  Control assets:
  Off-chain claims:

4. Actors
  Users:
  Operators:
  Privileged roles:
  External participants:
  Adversaries:

5. Entry points
  On-chain:
  Off-chain:
  Human/operational:

6. Trust boundaries
  Oracles:
  Bridges:
  RPC/indexers:
  Frontend/backend:
  Governance/admin:
  Custody/legal:

7. Assumptions
  Safety:
  Liveness:
  Economic:
  Operational:

8. Threat scenarios
  Threat:
  Preconditions:
  Attack path:
  Impact:
  Existing controls:
  Gaps:
  Priority:

9. Review plan
  Code areas:
  Invariants:
  Tests:
  Manual review questions:
  Monitoring requirements:

10. Residual risk
  Accepted:
  Mitigated:
  Requires decision:
\end{lstlisting}

This template is not bureaucracy. It is a compression format for
security reasoning.

\subsection{Anti-Patterns}

Avoid these.

{\def\LTcaptype{none} % do not increment counter
\begin{longtable}[]{@{}
  >{\raggedright\arraybackslash}p{(\linewidth - 2\tabcolsep) * \real{0.5000}}
  >{\raggedright\arraybackslash}p{(\linewidth - 2\tabcolsep) * \real{0.5000}}@{}}
\toprule\noalign{}
\begin{minipage}[b]{\linewidth}\raggedright
Anti-pattern
\end{minipage} & \begin{minipage}[b]{\linewidth}\raggedright
Why it fails
\end{minipage} \\
\midrule\noalign{}
\endhead
\bottomrule\noalign{}
\endlastfoot
Bug-list threat model & Names vulnerabilities without modeling
system-specific attack paths \\
Diagram-only threat model & Looks serious but does not drive tests or
review \\
Contract-only scope & Misses oracles, governance, frontends, keepers,
bridges, and keys \\
Honest-user model & Reviews product flows instead of adversarial
transitions \\
``Decentralized'' as assumption & Replaces precise validator,
governance, oracle, or admin assumptions with a slogan \\
Ignoring liveness & Preserves accounting while trapping funds \\
Treating trusted roles as safe & Hides the largest authority edges \\
No update process & Model becomes stale after deployment or governance
change \\
Equal-priority threats & Fails to guide scarce review effort \\
No residual-risk statement & Leaves users and teams unclear about what
remains trusted \\
\end{longtable}
}

The Threat Modeling Manifesto values finding and fixing design issues
over checkbox compliance, and doing threat modeling over talking about
it.\footnote{Threat Modeling Manifesto,
  https://www.threatmodelingmanifesto.org/.} That is the standard to use
here.

\subsection{What Good Looks Like}

A good threat model should let a reviewer say:

\begin{lstlisting}
I know where value is.
I know who can affect it.
I know which assumptions matter.
I know which assumptions are weakest.
I know which state transitions deserve the deepest review.
I know which invariants to write.
I know which tests should exist.
I know which external systems can break this protocol.
I know what risks the team is accepting.
\end{lstlisting}

If the threat model does not create that clarity, keep working.

\section{Worked Example: Threat Modeling a Lending Market}

Consider a simplified lending market:

\begin{lstlisting}
Users deposit collateral.
Users borrow stablecoin up to a collateral factor.
Oracle prices determine account health.
Liquidators repay debt and receive collateral at a discount.
Governance can list collateral assets and set parameters.
An admin multisig can pause markets and upgrade contracts after a timelock.
Keepers monitor accounts and submit liquidations.
\end{lstlisting}

\subsection{Step 1: Scope}

\begin{lstlisting}
Included:
  Lending market contracts
  Oracle integration
  Collateral token integration
  Liquidation logic
  Governance parameter changes
  Upgrade and pause controls

Excluded but documented:
  Frontend implementation
  Stablecoin issuer reserves
  External AMM liquidity sources
\end{lstlisting}

\subsection{Step 2: Assets}

{\def\LTcaptype{none} % do not increment counter
\begin{longtable}[]{@{}
  >{\raggedright\arraybackslash}p{(\linewidth - 4\tabcolsep) * \real{0.3333}}
  >{\raggedright\arraybackslash}p{(\linewidth - 4\tabcolsep) * \real{0.3333}}
  >{\raggedright\arraybackslash}p{(\linewidth - 4\tabcolsep) * \real{0.3333}}@{}}
\toprule\noalign{}
\begin{minipage}[b]{\linewidth}\raggedright
Asset
\end{minipage} & \begin{minipage}[b]{\linewidth}\raggedright
Representation
\end{minipage} & \begin{minipage}[b]{\linewidth}\raggedright
Failure
\end{minipage} \\
\midrule\noalign{}
\endhead
\bottomrule\noalign{}
\endlastfoot
Collateral & Token balances and account collateral records & Theft,
mispricing, frozen withdrawal \\
Debt & Borrow balances and interest indexes & Bad debt, incorrect
repayment \\
Stablecoin liquidity & Stablecoin balances and mint/transfer ability &
Insolvency, liquidity shortage \\
Liquidation right & Public liquidation function & Griefing, unfair
seizure, failed liquidation \\
Governance control & Voting and timelock state & Capture, unsafe
listing \\
Upgrade authority & Proxy admin/timelock & Malicious implementation \\
Oracle price & Oracle feed and cached values & Manipulated or stale
health checks \\
\end{longtable}
}

\subsection{Step 3: Actors}

{\def\LTcaptype{none} % do not increment counter
\begin{longtable}[]{@{}ll@{}}
\toprule\noalign{}
Actor & Capability \\
\midrule\noalign{}
\endhead
\bottomrule\noalign{}
\endlastfoot
Depositor & Supply collateral, withdraw \\
Borrower & Borrow, repay, increase leverage \\
Liquidator & Repay debt and seize collateral \\
MEV searcher & Order liquidation and oracle-update transactions \\
Governance & List assets, set factors, change oracle \\
Admin multisig & Pause and upgrade \\
Oracle operator & Publish price \\
Keeper & Monitor and call liquidations \\
Token issuer & Freeze or alter token behavior if token supports it \\
\end{longtable}
}

\subsection{Step 4: Entry Points}

\begin{lstlisting}
deposit(asset, amount)
withdraw(asset, amount)
borrow(asset, amount)
repay(asset, amount)
liquidate(account, repayAmount)
setCollateralFactor(asset, factor)
setOracle(asset, oracle)
pauseMarket(asset)
upgradeTo(implementation)
\end{lstlisting}

\subsection{Step 5: Assumptions}

{\def\LTcaptype{none} % do not increment counter
\begin{longtable}[]{@{}
  >{\raggedright\arraybackslash}p{(\linewidth - 4\tabcolsep) * \real{0.3333}}
  >{\raggedright\arraybackslash}p{(\linewidth - 4\tabcolsep) * \real{0.3333}}
  >{\raggedright\arraybackslash}p{(\linewidth - 4\tabcolsep) * \real{0.3333}}@{}}
\toprule\noalign{}
\begin{minipage}[b]{\linewidth}\raggedright
Assumption
\end{minipage} & \begin{minipage}[b]{\linewidth}\raggedright
Type
\end{minipage} & \begin{minipage}[b]{\linewidth}\raggedright
Failure
\end{minipage} \\
\midrule\noalign{}
\endhead
\bottomrule\noalign{}
\endlastfoot
Oracle price is fresh and manipulation-resistant & Safety &
Undercollateralized borrow \\
Liquidators act when profitable & Liveness & Bad debt grows \\
Collateral tokens behave as integrated & Safety & Accounting drift or
freeze \\
Governance cannot be cheaply captured & Safety & Malicious parameters \\
Pause cannot permanently trap users without documented authority &
Liveness & Frozen funds \\
Upgrade timelock gives users a real exit window & Safety/liveness &
Malicious upgrade before exit \\
\end{longtable}
}

\subsection{Step 6: Threat Scenarios}

Threat:

\begin{lstlisting}
Oracle manipulation enables bad debt.
\end{lstlisting}

Attack path:

\begin{lstlisting}
Attacker manipulates low-liquidity reference market.
Oracle reads manipulated price.
Attacker deposits inflated collateral.
Attacker borrows maximum stablecoin.
Price normalizes.
Account becomes insolvent.
\end{lstlisting}

Controls:

\begin{lstlisting}
TWAP or robust oracle
staleness checks
collateral caps
borrow caps
liquidation incentives
post-borrow health check
monitoring for price deviation
\end{lstlisting}

Threat:

\begin{lstlisting}
Governance lists unsafe collateral.
\end{lstlisting}

Attack path:

\begin{lstlisting}
Attacker accumulates or borrows voting power.
Proposal lists illiquid token with high collateral factor.
Attacker manipulates token price or mints supply.
Attacker borrows stablecoin and leaves bad debt.
\end{lstlisting}

Controls:

\begin{lstlisting}
listing framework
governance timelock
collateral caps
oracle review
liquidity requirements
risk committee veto
monitoring for proposal payloads
\end{lstlisting}

Threat:

\begin{lstlisting}
Liquidation liveness fails during market stress.
\end{lstlisting}

Attack path:

\begin{lstlisting}
Prices fall quickly.
Many accounts become unhealthy.
Liquidation transactions become expensive or unprofitable.
Keepers stop executing.
Bad debt accumulates.
\end{lstlisting}

Controls:

\begin{lstlisting}
liquidation incentives sized for stress
permissionless liquidation
batch limits
keeper redundancy
bad debt monitoring
emergency parameter controls
\end{lstlisting}

\subsection{Step 7: Review Plan}

Review priorities:

\begin{lstlisting}
1. Oracle integration and staleness/deviation handling.
2. Collateral and debt accounting invariants.
3. Liquidation math and liveness under stress.
4. Governance/admin authority and timelock bypasses.
5. Token integration assumptions.
6. Upgrade and pause behavior.
\end{lstlisting}

Invariants:

\begin{lstlisting}
total debt accounting matches user debts plus indexes
borrow cannot leave account below required health at execution time
liquidation cannot seize more collateral than rules allow
withdraw cannot make account unhealthy
collateral caps cannot be exceeded
paused state does not create undocumented permanent fund lock
\end{lstlisting}

Tests:

\begin{lstlisting}
manipulated oracle price
stale oracle price
fee-on-transfer collateral
collateral decimals mismatch
rapid price drop with many liquidations
governance parameter edge cases
upgrade queued while withdrawals are paused
\end{lstlisting}

This is a threat model that guides work. It tells the reviewer where to
look and what properties to test.

\section{Practical Discipline: The One-Page Threat Model}

Before a deep review, write one page with:

\begin{lstlisting}
Protocol goal:

Assets:

Critical transitions:

Privileged roles:

External dependencies:

Top safety assumptions:

Top liveness assumptions:

Most plausible profitable attack:

Worst privileged-role failure:

Most fragile dependency:

Top five review priorities:
\end{lstlisting}

If this page cannot be written, the protocol is not understood yet.

\section{Review Questions}

\begin{enumerate}
\def\labelenumi{\arabic{enumi}.}
\tightlist
\item
  What should a blockchain threat model produce besides a list of possible attackers?
\item
  Why should assets, actors, entry points, privileged roles, and trust boundaries be mapped before naming bug classes?
\item
  What is the difference between a user, attacker, rational participant, trusted operator, and privileged actor?
\item
  Why is an external dependency a security assumption rather than background infrastructure?
\item
  How does an attack tree help convert a vague threat into review priorities?
\item
  What makes a threat model useful during review rather than decorative?
\item
  Why should a threat model include both safety and liveness assumptions?
\item
  How should a reviewer turn threat-model output into tests, invariants, mitigations, and monitoring?
\end{enumerate}

\section{Applied Exercises}

\subsection{Exercise 1: One-Page Threat Model}

Given a protocol description, produce a one-page threat model with:

\begin{itemize}
\tightlist
\item
  protocol goal;
\item
  assets;
\item
  actors;
\item
  critical transitions;
\item
  privileged roles;
\item
  external dependencies;
\item
  top safety assumptions;
\item
  top liveness assumptions;
\item
  top five review priorities.
\end{itemize}

Expected output: a one-page artifact, not an essay.

\subsection{Exercise 2: Data and Control Flow}

Draw a data/control-flow diagram for one of the following:

\begin{itemize}
\tightlist
\item
  lending market with oracle and liquidators;
\item
  vault with strategist and withdrawal queue;
\item
  bridge with relayers and validator signatures.
\end{itemize}

For each arrow, state what crosses the boundary, who can modify it, who verifies it, and what happens if it is stale, false, missing, delayed, replayed, or reordered.

Expected output: one diagram and one trust-boundary table.

\subsection{Exercise 3: Attack Tree to Review Plan}

Start with this goal:

\begin{lstlisting}
Cause the protocol to lose assets or freeze user exits.
\end{lstlisting}

Build an attack tree with at least three branches. Then convert it into a review plan with depth budgets.

Expected output: attack tree plus a prioritized review plan.

\subsection{Exercise 4: Threat Model Critique}

Critique this weak statement:

\begin{lstlisting}
The system is safe because only trusted actors can perform sensitive actions.
\end{lstlisting}

Rewrite it as concrete assumptions about roles, keys, thresholds, timelocks, observability, and recovery paths.

Expected output: an assumption register and one residual-risk paragraph.

\section{Key Takeaways}

\begin{itemize}
\tightlist
\item
  Threat modeling turns a protocol description into a security model
  that guides review.
\item
  A useful threat model connects system scope, assets, actors, entry
  points, assumptions, threats, mitigations, and tests.
\item
  Blockchain threat models must include chain behavior, transaction
  ordering, economic incentives, external dependencies, infrastructure,
  and human operations.
\item
  Users are not trusted to follow product intent. Public entry points
  must survive adversarial but valid use.
\item
  Rational participants, Byzantine actors, compromised keys, insiders,
  and negligent operators require different controls.
\item
  Trust assumptions should be written in falsifiable form and classified
  by safety, liveness, or both.
\item
  Attack surfaces exist across chain, contract, economic, cross-domain,
  infrastructure, and human layers.
\item
  STRIDE and vulnerability lists are prompts, not substitutes for
  system-specific reasoning.
\item
  A good threat model produces review priorities, invariants, tests,
  mitigations, monitoring, and residual-risk decisions.
\item
  A threat model that does not change what gets reviewed or tested is
  decoration.
\end{itemize}

\section{Further Reading}

\begin{itemize}
\tightlist
\item
  OWASP, \href{https://owasp.org/www-community/Threat_Modeling}{``Threat
  Modeling''}. Good general overview of threat-modeling goals,
  lifecycle, and the four-question framework.
\item
  \href{https://www.threatmodelingmanifesto.org/}{``Threat Modeling
  Manifesto''}. Concise principles and anti-patterns for practical
  threat modeling.
\item
  Microsoft,
  \href{https://learn.microsoft.com/en-us/azure/security/develop/threat-modeling-tool-threats}{``Microsoft
  Threat Modeling Tool threats''}. Practical overview of STRIDE threat
  categories.
\item
  NIST,
  \href{https://nvlpubs.nist.gov/nistpubs/Legacy/SP/nistspecialpublication800-30r1.pdf}{``SP
  800-30 Revision 1: Guide for Conducting Risk Assessments''}. Useful
  for threat sources, likelihood, impact, and risk framing.
\item
  OWASP,
  \href{https://owasp.org/www-project-smart-contract-top-10/}{``Smart
  Contract Top 10''}. Blockchain-specific awareness list of common
  smart-contract risk categories.
\item
  OWASP Smart Contract Security,
  \href{https://scs.owasp.org/SCSVS/controls/SCSVS-ARCH-3/}{``S1.3
  Threat Modeling''}. Smart-contract security verification requirements
  for identifying threats, assessing risks, and implementing
  mitigations.
\item
  OWASP Smart Contract Security,
  \href{https://scs.owasp.org/sctop10/Web3-Attack-Vectors-Top15/}{``Alternate
  Top 15 - Web3 Attack Vectors Beyond Smart Contracts''}. Useful
  reminder that protocol risk includes infrastructure, keys, supply
  chain, and social engineering.
\end{itemize}
